\subsection{相干光照下，电子何时仍可视为纯态}
\label{sec:electron-quantum-exchange-coherence}

经典 PINEM 给出一个纯电子态。若把同一束相干光完全量子化，电子输出是否仍等于这个纯态，需要对光场求偏迹才能判断。相干态满足 $\hat a|\alpha_{\rm coh}\rangle=\alpha_{\rm coh}|\alpha_{\rm coh}\rangle$，吸收项因此产生熟悉的经典振幅；发射项含有 $\hat a^\dagger$，还会产生不能用这个数代替的分量。下面保留两项，计算电子的密度矩阵、能谱和纯度。光子统计在电子谱中的表现由同一联合态描述。\cite{Dahan2021Science}

在已建立的单模、双边平移、低反冲交换模型中，去掉不影响观测的整体 Magnus 相位，令
\begin{equation*}
 S_Q=\exp(\gq\hat b^\dagger\hat a-\gq^*\hat b\hat a^\dagger),
 \qquad r_Q=|\gq|^2,\qquad \beta=\alpha_{\rm coh}\gq.
\end{equation*}
$r_Q$ 是无量纲单电子真空交换强度；它不同于由宏观场增强的 $|\beta|^2$。对电子中心通道 $|0\rangle_e$ 和相干光输入，定义经典 PINEM 纯态
\begin{equation*}
 |\psi_\beta\rangle=S_{\rm tr}(\beta)|0\rangle
 =\sum_\ell J_\ell(2|\beta|)\ee^{\ii\ell\arg\beta}|\ell\rangle.
\end{equation*}
精确约化电子态可写成一组向损失方向平移的纯态混合：
\begin{equation}
 \rho_e^{\rm out}=\sum_{m=0}^{\infty}w_m\,
 \hat b^m|\psi_\beta\rangle\langle\psi_\beta|(\hat b^\dagger)^m,
 \qquad w_m=\ee^{-r_Q}\frac{r_Q^m}{m!}.
 \label{eq:electron-coherent-light-reduced-state}
\end{equation}
每个 $\hat b^m|\psi_\beta\rangle$ 都是把经典输出向能量损失方向移动 $m$ 阶后的归一态，$w_m$ 是它的概率权重。局部推导中会把光的平均相干幅度移去，余下的不同光子数对应这些不同的电子平移。只观测电子时，对正交的光子末态求和，便得到上式的混合。

能谱、平均能量交换和方差随即成为
\begin{equation}
 \begin{aligned}
 P_\ell&=\sum_{m=0}^{\infty}\ee^{-r_Q}\frac{r_Q^m}{m!}
 J_{\ell+m}^2(2|\beta|),\\
 \langle\ell\rangle&=-r_Q,\qquad
 \operatorname{Var}\ell=2|\beta|^2+r_Q.
 \end{aligned}
 \label{eq:electron-coherent-light-spectrum-moments}
\end{equation}
负的均值表示电子向光场净提供能量 $r_Q\hbar\omega$。经典场贡献给出对称展宽，真空项给出额外的损失偏置；这两项在一个概率卷积中同时出现。

为了量化混合程度，使用\emph{纯度} $\operatorname{Tr}\rho^2$。若归一化密度矩阵的本征值为 $w_m$，纯度就是 $\sum_mw_m^2$；只有一个本征值等于一时为纯态，几个非零权重共同存在时纯度小于一。这里的平移态彼此正交：$S_{\rm tr}$ 与 $\hat b$ 对易，而 $\langle0|\hat b^{m-n}|0\rangle=\delta_{mn}$。因此已经找到了本征值 $w_m$，无需另求矩阵对角化，得到
\begin{equation}
 \operatorname{Tr}[(\rho_e^{\rm out})^2]
 =\sum_{m=0}^{\infty}w_m^2
 =\ee^{-2r_Q}I_0(2r_Q),
 \label{eq:electron-coherent-light-purity}
\end{equation}
其中 $I_0(x)=\sum_{m\ge0}(x^2/4)^m/(m!)^2$ 为零阶第一类修正 Bessel 函数。它与振荡型的 $J_0$ 不同。这个纯度在本模型中只依赖 $r_Q$；增大相干光幅度会增加经典调制，却不会自动消除电子留下的真空量子记录。

\begin{exbox}[经典纯态近似的误差怎样随耦合缩小]
以 $|\psi_\beta\rangle$ 为经典参考态，求量子电子态偏离它的大小，并比较“能谱接近”与“量子态接近”两种判据。

前面已证明各平移态彼此正交；未平移项的本征值为 $w_0=\ee^{-r_Q}$。因此与参考纯态的重叠概率为
\begin{equation*}
 F_\beta\equiv\langle\psi_\beta|\rho_e^{\rm out}|\psi_\beta\rangle
 =\ee^{-r_Q}.
\end{equation*}
采用\emph{迹距离} $\mathcal D_{\rm tr}(\rho,\sigma)=\tfrac12\|\rho-\sigma\|_1$，其中迹范数是奇异值之和；对厄米差矩阵，就是本征值绝对值之和。这里差矩阵的本征值是 $w_0-1,w_1,w_2,\ldots$，故无需做无限维对角化：
\begin{equation*}
 \mathcal D_{\rm tr}
 =\frac12\bigl[(1-w_0)+\sum_{m\ge1}w_m\bigr]
 =1-\ee^{-r_Q}.
\end{equation*}
迹距离控制任意一个测量事件的概率偏差。要求该误差不超过 $0<\varepsilon<1$，等价于
\begin{equation*}
 r_Q\le-\ln(1-\varepsilon).
\end{equation*}
小耦合时，$\mathcal D_{\rm tr}=r_Q+O(r_Q^2)$，纯度则为
$\operatorname{Tr}[(\rho_e^{\rm out})^2]=1-2r_Q+3r_Q^2+O(r_Q^3)$。

相反，量子方差相对于经典方差的增量在 $\beta\ne0$ 时为
\begin{equation*}
 \frac{\operatorname{Var}\ell-2|\beta|^2}{2|\beta|^2}
 =\frac{r_Q}{2|\beta|^2}
 =\frac1{2|\alpha_{\rm coh}|^2}.
\end{equation*}
增大平均光子数会让这个\emph{相对方差误差}缩小，却不改变固定 $r_Q$ 下的迹距离。要同时恢复完整的经典电子纯态，应取 $\gq\to0$、$|\alpha_{\rm coh}|\to\infty$ 且 $\beta$ 固定。在这个极限中，经典调制保持不变，而真空交换权重趋于零。
\end{exbox}

弱 Fock 交换还需要单独检查玻色增强。对确定光子数 $n$，截断单量子概率需要 $r_Q(n+1)\ll1$；对有宽尾的光子分布，应检查相应高阶光子矩及实际占据尾。于是 $P_{+1}\simeq r_Q\langle n\rangle$ 与高阶阶乘矩公式都不是只由 $r_Q\ll1$ 保证。经典场和量子场两种近似条件应分别报告。

\begin{derivation}[移去光场平均值后，剩下真空记录]
用位移算符 $D(\alpha)=\exp(\alpha\hat a^\dagger-\alpha^*\hat a)$，有 $|\alpha\rangle=D(\alpha)|0\rangle$ 与 $D^\dagger\hat aD=\hat a+\alpha$。因此
\begin{equation*}
 D^\dagger S_QD
 =\exp[\gq\hat b^\dagger\hat a-\gq^*\hat b\hat a^\dagger
       +\beta\hat b^\dagger-\beta^*\hat b]
 =S_Q S_{\rm tr}(\beta).
\end{equation*}
最后一步使用双边平移的 $[\hat b,\hat b^\dagger]=0$，使量子部分与经典部分对易。对光真空作用时，令 $A=-\gq^*\hat b\hat a^\dagger$、$B=\gq\hat b^\dagger\hat a$，则 $[A,B]=r_Q I$，正规排序给
\begin{equation*}
 S_Q|0,0\rangle=\ee^{-r_Q/2}\sum_{m=0}^{\infty}
 \frac{(-\gq^*)^m}{\sqrt{m!}}|-m,m\rangle.
\end{equation*}
移回的 $D(\alpha)$ 只作用于被取迹的光场，不改变电子约化态；不同 $m$ 的光记录正交，因而得到\eqnrefs{eq:electron-coherent-light-reduced-state}。谱中随机整数可写作 $\ell=L_\beta-M$，其中 $P(L_\beta=j)=J_j^2(2|\beta|)$、$M$ 服从均值 $r_Q$ 的 Poisson 分布且二者独立。利用 $\langle L_\beta\rangle=0$、$\operatorname{Var}L_\beta=2|\beta|^2$，就得到\eqnrefs{eq:electron-coherent-light-spectrum-moments}。
\end{derivation}
